\documentclass[11pt]{article}
\usepackage[margin=1in]{geometry}
\usepackage{amsmath,amssymb,amsthm}
\usepackage{algorithm}
\usepackage{algpseudocode}
\usepackage{booktabs}
\usepackage{hyperref}

\newtheorem{definition}{Definition}
\newtheorem{theorem}{Theorem}
\newtheorem{lemma}[theorem]{Lemma}
\newtheorem{property}{Property}

\title{\textbf{Mathematical Specification, Numerical Stabilization, and Positive Weight Dynamics of Binary Cross-Entropy with Logits (ALGO-NN-15)}}
\author{Team Scaibu \\ \small Email: \texttt{jaydeep.9664920749@gmail.com} \\ \small Architecture and Core Engine Team}
\date{\today}

\begin{document}
\maketitle

\begin{abstract}
Binary classification and multi-label classification require evaluating log-odds through the logistic loss. Direct composition of the sigmoid function $\sigma(z)$ with logarithm operations suffers from underflow and infinite gradients. This specification details the log-sum-exp stabilization of Binary Cross-Entropy with Logits ($\mathcal{L}_{\mathrm{BCE}}$), incorporates positive class reweighting $p_w$, proves gradient boundedness, derives step-by-step complexity, and provides full pseudocode verified against the reference implementation.
\end{abstract}

\section{Notation and Mathematical Preliminaries}
\label{sec:notation}

Let $z \in \mathbb{R}$ represent an unnormalized logit (log-odds $\log\frac{p}{1-p}$), and $y \in [0, 1]$ denote the ground-truth target probability. Let $p_w \in \mathbb{R}_{>0}$ denote the positive class weighting multiplier.

\subsection{Coordinate and Sigmoidal Mechanics}
\paragraph{Intuitive Meaning:} The sigmoid function maps unconstrained real numbers to probabilities in $(0, 1)$. Naively calculating $\log(\sigma(z))$ for large negative $z$ causes floating-point underflow to zero, producing $\log(0) = -\infty$. Stabilized formulation combines the logarithm and exponential into a single branch.

\paragraph{Formal Mathematical Definition:}
\begin{equation}
\sigma(z) = \frac{1}{1 + e^{-z}} = \frac{e^z}{1 + e^z}.
\end{equation}
The stable binary cross-entropy with positive weighting $p_w$ is:
\begin{equation}
\label{eq:bce_stable}
\mathcal{L}(z, y; p_w) = \begin{cases}
(1 - y) z + (1 + (p_w - 1)y) \log(1 + e^{-z}) & \text{if } z \ge 0, \\
(1 + (p_w - 1)y) \log(1 + e^z) - p_w y z & \text{if } z < 0.
\end{cases}
\end{equation}

\paragraph{Worked Numerical Example:}
Let $z = 2.0, y = 1.0, p_w = 1.0$:
$z \ge 0 \implies \mathcal{L} = (0)(2.0) + (1.0)\log(1 + e^{-2.0}) = \log(1 + 0.135335) = \log(1.135335) = 0.126928$.
Gradient: $\sigma(2.0) - y = 0.880797 - 1.0 = -0.119203$.

\subsection{Summary of Symbols}
\begin{itemize}
    \item $N \in \mathbb{N}$: Number of instances in the batch.
    \item $\mathbf{z} \in \mathbb{R}^N$: Logit prediction vector.
    \item $\mathbf{y} \in [0, 1]^N$: Target probability label vector.
    \item $p_w \in \mathbb{R}_{>0}$: Positive class loss weighting scalar ($\mathrm{pos\_weight}$).
    \item $\mathbf{w} \in \mathbb{R}_{\ge 0}^N$: Per-element weight vector.
    \item $\boldsymbol{\sigma} \in (0, 1)^N$: Predicted sigmoid probability vector.
    \item $\mathbf{g} = \nabla_{\mathbf{z}} \mathcal{L} \in \mathbb{R}^N$: Analytic gradient vector.
\end{itemize}

\section{Problem Definition and Core Concepts}
\label{sec:problem_definition}

\subsection{Positive Class Weighting Mechanics}
\paragraph{Intuition:} In extreme imbalance (e.g., $1\%$ positive labels), standard loss is minimized by predicting all negatives. Setting $p_w = 99.0$ scales the loss penalty on positive targets, forcing the optimizer to balance recall and precision.

\paragraph{Formal Gradient Derivation:}
Differentiating Equation~\eqref{eq:bce_stable} with respect to $z$:
\begin{equation}
\frac{\partial \mathcal{L}}{\partial z} = \sigma(z) \cdot \left[ 1 + (p_w - 1)y \right] - p_w y.
\end{equation}

\section{Algorithmic Specification}
\label{sec:algorithm}

Algorithm~\ref{alg:bce_logits} details the exact stable calculation procedure matching \texttt{impl.py}.

\begin{algorithm}[H]
\caption{Stable Binary Cross-Entropy with Logits and Positive Class Weighting}
\label{alg:bce_logits}
\begin{algorithmic}[1]
\Require Logits $\mathbf{z} \in \mathbb{R}^N$, targets $\mathbf{y} \in [0, 1]^N$, $p_w > 0$, weights $\mathbf{w} \in \mathbb{R}^N$, reduction $\mathrm{red} \in \{\text{mean}, \text{sum}, \text{none}\}$.
\Ensure Loss $\mathcal{L}$, probabilities $\boldsymbol{\sigma} \in \mathbb{R}^N$, gradients $\mathbf{g} \in \mathbb{R}^N$.
\State Initialize $\mathbf{l} \gets [], \boldsymbol{\sigma} \gets [], \mathbf{g} \gets []$
\For{$i = 1$ to $N$}
    \State $z \gets z_i, \quad y \gets y_i, \quad w \gets w_i$ (or $1.0$)
    \If{$z \ge 0$}
        \State $e_{-z} \gets \exp(-z)$
        \State $\sigma_i \gets \frac{1.0}{1.0 + e_{-z}}$
        \State $\mathrm{log\_term} \gets \log(1.0 + e_{-z})$
        \State $l_i \gets (1.0 - y) \cdot z + (1.0 + (p_w - 1.0) \cdot y) \cdot \mathrm{log\_term}$
    \Else
        \State $e_z \gets \exp(z)$
        \State $\sigma_i \gets \frac{e_z}{1.0 + e_z}$
        \State $\mathrm{log\_term} \gets \log(1.0 + e_z)$
        \State $l_i \gets (1.0 + (p_w - 1.0) \cdot y) \cdot \mathrm{log\_term} - p_w \cdot y \cdot z$
    \EndIf
    \State $l_i \gets w \cdot l_i$
    \State $g_i \gets w \cdot \left[ \sigma_i \cdot (1.0 + (p_w - 1.0) \cdot y) - p_w \cdot y \right]$
    \State Append $l_i$ to $\mathbf{l}$, append $\sigma_i$ to $\boldsymbol{\sigma}$, append $g_i$ to $\mathbf{g}$
\EndFor
\If{$\mathrm{red} = \text{mean}$}
    \State $W_{\mathrm{norm}} \gets \max(\sum w_i, 1.0)$
    \State \Return $\frac{1}{W_{\mathrm{norm}}}\sum l_i, \quad \boldsymbol{\sigma}, \quad \frac{1}{W_{\mathrm{norm}}}\mathbf{g}$
\ElsIf{$\mathrm{red} = \text{sum}$}
    \State \Return $\sum l_i, \quad \boldsymbol{\sigma}, \quad \mathbf{g}$
\Else
    \State \Return $\mathbf{l}, \quad \boldsymbol{\sigma}, \quad \mathbf{g}$
\EndIf
\end{algorithmic}
\end{algorithm}

\section{Theoretical Guarantees}
\label{sec:guarantees}

\begin{theorem}[Global Gradient Boundedness and Saturation Limits]
\label{thm:bce_grad_bound_full}
\textbf{Assumptions:} Let target $y \in [0, 1]$, positive weight $p_w > 0$, and logit $z \in \mathbb{R}$.
\textbf{Guarantee:} The gradient $\frac{\partial \mathcal{L}}{\partial z}$ is strictly bounded in the range $[-p_w, 1]$ for all $z \in \mathbb{R}$.
\textbf{Proof:}
Since $\sigma(z) \in (0, 1)$, we evaluate the extrema:
\begin{equation}
\inf_{z \in \mathbb{R}} \frac{\partial \mathcal{L}}{\partial z} = 0 \cdot (1 + (p_w - 1)y) - p_w y = -p_w y \ge -p_w.
\end{equation}
\begin{equation}
\sup_{z \in \mathbb{R}} \frac{\partial \mathcal{L}}{\partial z} = 1 \cdot (1 + (p_w - 1)y) - p_w y = 1 - y \le 1.
\end{equation}
Because the derivative is continuous and monotonic in $z$, the gradient is strictly bounded within $[-p_w, 1]$.
\textbf{Limitations:} When $|z| \to \infty$ with matching label, gradient vanishes $\to 0$ (saturation), which can slow learning if initialization is poorly scaled.
\end{theorem}

\section{Complexity Analysis}
\label{sec:complexity}

\begin{itemize}
    \item \textbf{Time Complexity:} $\mathcal{O}(N)$ elementary exponential, logarithm, and arithmetic operations.
    \item \textbf{Space Complexity:} $\mathcal{O}(N)$ memory for probability and gradient arrays.
\end{itemize}

\section{Worked Numerical Example}
\label{sec:worked_example}

Let $z = -1.0, y = 1.0, p_w = 2.0$:
\begin{enumerate}
    \item $z < 0 \implies e_z = \exp(-1.0) = 0.367879$.
    \item Probability: $\sigma = \frac{0.367879}{1.367879} = 0.268941$.
    \item $\log(1 + e^z) = \log(1.367879) = 0.313262$.
    \item Loss: $\mathcal{L} = (1 + (2 - 1)(1))(0.313262) - (2)(1)(-1.0) = 2(0.313262) + 2.0 = 2.626524$.
    \item Gradient: $g = 0.268941(1 + 1) - 2.0(1.0) = 2(0.268941) - 2.0 = 0.537882 - 2.0 = -1.462118$.
\end{enumerate}

\section{Numerical Considerations and Edge Cases}
\label{sec:numerical}

\begin{itemize}
    \item \textbf{Branching by Sign:} Splitting computation at $z \ge 0$ vs $z < 0$ ensures that $\exp(-|z|) \le 1.0$, completely avoiding floating-point overflow for any finite float.
    \item \textbf{Soft Targets:} Continuous targets $y \in (0, 1)$ (e.g. from mixup or distillation) are fully supported by the linear convex combination.
\end{itemize}

\begin{thebibliography}{9}
\bibitem{bishop2006} C.~M.~Bishop, \emph{Pattern Recognition and Machine Learning}, Springer, 2006.
\end{thebibliography}

\end{document}
